\chapter{Présentation des sources et audit de qualité}
\label{chap:audit}

Les notebooks 01 et 02 examinent les fichiers avant la préparation. Le premier permet
de comprendre les relations sur 2024~; le second étend les contrôles aux six années.
Leurs effectifs concernent les sources complètes, métropole et outre-mer. Les résultats
du filtre métropolitain seront présentés au chapitre~\ref{chap:preparation}.

\section{Sources et documentation de référence}
Les extractions annuelles BAAC sont publiées sur data.gouv.fr~\cite{baac}. Le document
de codification utilisé est la \textit{Description des bases de données annuelles des
accidents corporels de la circulation routière, années 2005 à 2024}, datée du
21 octobre 2025~\cite{doc-onisr}. Les constats du rapport correspondent aux 24 fichiers
CSV fournis pour le projet et vérifiés le 16 octobre 2026. Une publication ultérieure
pourrait modifier ces effectifs~; les empreintes des fichiers permettent d'identifier
la version réellement analysée.

La lecture utilise le séparateur point-virgule et préserve les identifiants en texte.
Un code numérique ne devient pas pour autant une mesure~: \code{lum}, \code{catv}
et \code{grav} désignent des catégories. Calculer une moyenne de leurs codes leur
attribuerait une signification quantitative qui n'est pas celle de la documentation.

\section{Quatre fichiers et plusieurs niveaux d'observation}
\begin{table}[htbp]\centering\small\renewcommand{\arraystretch}{1.2}
\caption{Rôle des fichiers annuels.}\label{tab:tables-source}
\begin{tabularx}{\textwidth}{@{}lLL@{}}\toprule
Fichier & Une ligne représente & Informations principales\\\midrule
\code{caracteristiques} & Un accident & Date, heure, lumière, météo, localisation.\\
\code{vehicules} & Un véhicule impliqué & Catégorie, manœuvre, obstacles, choc.\\
\code{usagers} & Un usager impliqué & Catégorie, naissance, gravité, équipements.\\
\code{lieux} & Un enregistrement de lieu & Route, circulation, profil, surface, limitation.\\\bottomrule
\end{tabularx}
\sourcecalcul{documentation ONISR et contrôles des notebooks 01 et 02}
\end{table}

La figure~\ref{fig:relations} représente les liens contrôlés au sein d'un millésime.
Le couple accident-véhicule permet de rattacher les usagers au véhicule correspondant.
Un conducteur et son passager partagent ce véhicule, tout en restant deux observations.
La documentation précise qu'un piéton est rattaché au véhicule qui l'a heurté~:
\code{catv} décrit alors ce véhicule et non le mode de déplacement du piéton.

\begin{figure}[htbp]\centering
\begin{tikzpicture}[box/.style={draw=anthracite!60,rounded corners=2pt,fill=black!2,text width=5.1cm,minimum height=1.35cm,align=left,font=\small},
  edge/.style={-{Stealth[length=2.1mm]},thick,draw=terreCuite},every edge quotes/.style={font=\scriptsize}]
\node[box] (c) at (-3.7,0) {\textbf{Caractéristiques}\\Une ligne par accident\\Clé : \code{Num_Acc}};
\node[box] (v) at (3.7,0) {\textbf{Véhicules}\\Une ligne par véhicule impliqué\\Clé : accident + \code{id_vehicule}};
\node[box] (l) at (-3.7,-3.2) {\textbf{Lieux sources}\\Une ou plusieurs lignes par accident\\Aucune clé véhicule dans le CSV};
\node[box] (u) at (3.7,-3.2) {\textbf{Usagers}\\Une ligne par usager source\\Clé technique créée pour 2019-2024};
\draw[edge] (c.east) -- node[above,font=\scriptsize,align=center]{1 vers plusieurs\\\code{Num_Acc}} (v.west);
\draw[edge] (c.south) -- node[left,font=\scriptsize,align=right]{1 vers 1 à 5\\lignes observées} (l.north);
\draw[edge] (v.south) -- node[left,font=\scriptsize,align=right]{1 vers 0 à plusieurs\\accident + véhicule} (u.north);
\draw[edge] (c.south east) -- node[below,sloped,font=\scriptsize,align=center]{1 vers plusieurs\\\code{Num_Acc}} (u.north west);
\end{tikzpicture}
\caption{Relations contrôlées entre les tables sources. Les cardinalités décrivent les fichiers examinés ; elles sont vérifiées avant chaque jointure.}
\label{fig:relations}
\sourcecalcul{schéma du projet, à partir des notebooks 01 et 02}
\end{figure}

\section{Dimensions et différences de schéma}
\begin{table}[htbp]\centering\small
\caption{Nombre de lignes dans les sources complètes.}\label{tab:dimensions}
\begin{tabularx}{\textwidth}{@{}lLLLL@{}}\toprule
Année & Caractéristiques & Lieux & Véhicules & Usagers\\\midrule
2019 & 58\,840 & 58\,840 & 100\,710 & 132\,977\\
2020 & 47\,744 & 47\,744 & 81\,066 & 105\,295\\
2021 & 56\,518 & 56\,518 & 97\,315 & 129\,248\\
2022 & 55\,302 & 55\,302 & 94\,493 & 126\,662\\
2023 & 54\,822 & 70\,860 & 93\,585 & 125\,789\\
2024 & 54\,402 & 70\,248 & 92\,678 & 125\,187\\\bottomrule
\end{tabularx}
\sourcecalcul{notebook 02 ; nombre de lignes et non nombre de personnes pour toutes les colonnes}
\end{table}

Les caractéristiques comportent 15 colonnes, les lieux 18 et les véhicules 11. Les
usagers passent de 15 à 16 colonnes en 2021 avec l'apparition de \code{id_usager}.
En 2022, la clé des caractéristiques est nommée \code{Accident_Id}. Une concaténation
sans harmonisation ferait donc apparaître deux colonnes pour un même rôle.

Les espaces, dont des espaces insécables internes, concernent les identifiants de
véhicule et d'usager. La préparation les retirera dans toutes les tables liées.
Le contrôle vérifie qu'aucune valeur normalisée ne réunit plusieurs écritures brutes
distinctes dans une même colonne~: aucune collision n'est détectée.

\section{Clés, couverture et répétitions}
La clé accident est complète et unique dans les caractéristiques. La clé
\code{Num_Acc} + \code{id_vehicule} est unique dans les véhicules. Les identifiants
d'usager sont complets et uniques lorsqu'ils existent. Les ensembles d'accidents
des quatre tables sont identiques chaque année et chaque usager retrouve son véhicule.
Ces contrôles portent sur les identifiants eux-mêmes, pas seulement sur des effectifs
égaux.

La relation inverse est différente~: certains véhicules n'ont pas d'usager associé.
Sur les sources complètes, les nombres sont respectivement 2\,040, 1\,899, 6, 0, 40 et
24 entre 2019 et 2024. Le PDF explique les 24 cas de 2024 par des conducteurs en fuite
non créés dans la table des usagers~\cite{doc-onisr}. Cette explication ne doit pas être
étendue automatiquement à tous les écarts des autres années. Aucun usager ni aucune
gravité ne sont inventés pour compléter ces véhicules.

En 2019 et 2020, 101 et 63 lignes d'usagers supplémentaires sont identiques sur tous
leurs attributs sources. Sans identifiant individuel, ce constat ne prouve pas un
double enregistrement~: deux passagers peuvent avoir les mêmes caractéristiques.
Les lignes sont donc conservées et identifiées par leur rang source.

\section{Pourquoi les lieux demandent un traitement particulier}
Jusqu'en 2022, chaque accident possède une ligne de lieux. En 2023, 15\,870 accidents
ont plusieurs lignes, contre 15\,645 en 2024, soit respectivement 28,95\,\% et 28,76\,\%
des accidents des sources complètes. Le maximum observé est de cinq lignes. Certains
attributs varient et ces lignes ne sont donc pas de simples copies exactes.

L'accident \code{202400000005} illustre le problème. Quatre usagers sont associés à
deux lignes de lieux portant le même nom de voie, mais avec des limitations de 50 et
70 km/h et des profils de route différents. Une jointure brute sur \code{Num_Acc}
produirait huit lignes. La moyenne de 60 km/h, ou le choix de la première ligne,
n'aurait pas de justification comme limitation du lieu de chaque usager.

Dans sa réponse du 12 décembre 2024 à la discussion «~lieux-2023~», l'ONISR explique
que plusieurs lieux peuvent être décrits aux intersections, en lien avec les véhicules
concernés~\cite{discussion-lieux}. Les CSV reçus ne fournissent toutefois aucune clé
véhicule dans \code{lieux}. Le projet retiendra donc des informations communes au
niveau de l'accident. Cette explication ne démontre pas pourquoi la représentation
multiple apparaît à partir de 2023 dans ces exports.

\section{Codes, informations inconnues et plausibilité}
Le sens d'un code dépend de sa variable. \code{trajet = 0} désigne une information
non renseignée, alors que \code{secu1 = 0} signifie aucun équipement et
\code{vosp = 0} sans objet. Un remplacement global des zéros détruirait ces distinctions.
Certains codes observés sont absents de la rubrique du PDF, notamment 7 et 8 dans
\code{actp}. Ils sont signalés sans leur attribuer un sens supposé.

Les sources complètes contiennent 419 gravités inconnues~: 60 en 2021, 241 en 2022 et
118 en 2023. L'audit relève aussi 191 âges approximatifs supérieurs à 110 ans, issus
des naissances 1900 ou 1901 en 2019-2020. Des naissances manquantes apparaissent à partir
de 2021. Le PDF signale l'ajout d'usagers en fuite depuis 2021, avec des informations
parfois manquantes~\cite{doc-onisr}. Cette évolution constitue un contexte de collecte~;
elle ne permet pas d'identifier la cause de chaque valeur absente.

Les dates, heures et coordonnées passent les contrôles élémentaires appliqués~:
interprétation, concordance d'année et limites mondiales. Cela ne garantit pas que la
position GPS correspond au lieu réel. Le contrôle des communes métropolitaines révèle
trois discordances avec le département en 2019 et un code \code{N/C} en 2022. Sans
référentiel historique ni vérification complémentaire, l'audit ne choisit pas lequel
des deux renseignements corriger.

\FloatBarrier
\bilan{Les six années peuvent être réunies après harmonisation. Les risques principaux
concernent la multiplication des usagers par les lieux, la confusion entre codes
spéciaux et valeurs ordinaires, et une suppression injustifiée d'observations similaires.
Le chapitre suivant transforme ces constats en règles vérifiables.}
